\section{Introduction}

State-of-the-art (SOTA) deep learning models all share a common characteristic: they all have an extremely large number of parameters (10's if not 100's of billions parameters). Currently, only a few industry labs can pretrain large language models due to their high training cost. However, many pretrained models are accessible either through an API (GPT4, \cite{openai2023gpt4}) or through open-source platforms (Llama, \cite{touvron2023llama}). Most practitioners are interested in using such models for specific tasks and want to \emph{adapt} these models to a new, generally smaller task. This procedure is known as \emph{finetuning}, where one adjusts the weights of the pretrained model to improve performance on the new task. However, due to the size of SOTA models, adapting to down-stream tasks with full finetuning (finetuning all model parameters) is computationally infeasible as it requires modifying the weights of the pretrained models using gradient methods which is a costly process.  Besides, a model that has already learned generally useful representations during pretraining would not require  in-principle significant adaptation of all parameters. With this intuition, researchers have proposed a variety of resource-efficient finetuning methods which typically freeze the pretrained weights and tune only a small set of newly inserted parameters. Such methods include prompt tuning \citep{lester2021power} where a ``soft prompt" is learned and appended to the input, the adapters method \citep{houlsby2019parameter} where  lightweight ``adapter" layers are inserted and trained, and $(IA)^3$ \citep{liu2022few} where activation vectors are modified with learned scalings. Another resource-efficient method is known as \emph{Low Rank Adaptation} \citep{hu2021lora}, or simply LoRA. In LoRA finetuning, only a low rank matrix, called an \textit{adapter}, that is added to the pretrained weights is trainable. The training can be done with any optimizer and in practice a common choice is Adam \citep{kingma2014adam}. Since the trained adapter is low-rank, this effectively reduces the number of trainable parameters in the fine-tuning process, significantly decreasing the training cost. On many tasks such as instruction finetuning, LoRA has been shown to achieve comparable or better performance compared with full-finetuning \citep{wang2023far, liu2023improved}, although on complicated, long form generation tasks, it is not always as performant. The impressive performance and the computational savings of LoRA have contributed to it becoming an industry standard finetuning method. 

Efficient use of LoRA requires a careful choice of hyperparameters: the rank and the learning rate. While some theoretical guidelines on the choice of the rank in LoRA exist in the literature (see e.g. \citet{zeng2023expressive}), there are no principled guidelines on how to set the learning rate, apart from common choices  of order $1e$-$4$.
\begin{figure}
    \centering
    \includegraphics[width=.9\linewidth]{figures/intro_table.png}
    \caption{\small{The key difference between standard LoRA and LoRA$+$ is in how learning rates are set (the matrices $G_A$ and $G_B$ are `effective' gradients from AdamW) With standard LoRA, the learning rate is the same for $A$ and $B$, which provably leads to suboptimal learning when embedding dimension is large. In LoRA$+$, we set the learning rate of $B$ to be $\lambda \times$ that of $A$, where $\lambda \gg 1$ is fixed. We later provide guidelines on how to set $\lambda$.}}
    \vspace{-0.5cm}
    \label{fig:intro_table}
\end{figure}
% \begin{wrapfigure}{r}{0.55\textwidth}
%   \begin{center}
%        \includegraphics[width=.99\linewidth]{figures/intro_table.png}
%   \end{center}
%       \caption{\small{The key difference between standard LoRA and LoRA$+$ is in how learning rates are set (the matrices $G_A$ and $G_B$ are `effective' gradients after Adam processing) With standard LoRA, the learning rate is the same for $A$ and $B$, which provably leads to inefficient finetuning when embedding dimension is large. In LoRA$+$, we set the learning rate of $B$ to be $\lambda \times$ that of $A$, where $\lambda \gg 1$ is a fixed ratio. We later provide guidelines on how to choose $\lambda$.}}
%     \label{fig:intro_table}
% \end{wrapfigure}
\vspace{-0.3cm}
\paragraph{Related Work.} \citet{dettmers2023qlora} introduced a quantized version of LoRA (or QLoRA), which further reduces computation costs by quantizing pretrained weights down to as few as four bits. Using QLoRA enables fine-tuning Llama-65b \citep{touvron2023llama}, on a single consumer GPU while achieving competitive performance with full-finetuning. To further improve LoRA training with quantization, \citet{li2023loftq} introduced a new method called LoftQ for computing a better initialization for quantized training. Additional variations of LoRA have been proposed such as VeRA \citep{kopiczko2023vera} which freezes random weight tied adapters and learns vector scalings of the internal adapter activations. This achieves a further reduction in the number of trainable parameters while achieving comparable performance to LoRA on several NLP finetuning tasks. However, to the best of our knowledge, there is no principled guidance for setting LoRA learning rate which is the focus of our work. 
\vspace{-0.5cm}
\paragraph{Contributions.} We provide guidelines for setting the learning rate through a theory of scaling for neural networks. There is a significant number of works on the scaling of neural networks from the infinite width/depth perspective. The approach is simple: take the width/depth of a neural network to infinity,\footnote{Depending on the model, one might want to scale width with fixed depth and vice-versa, or both at the same time. See \cref{sec:scaling} for more details.} understand how the limit depends on the choice of the hyperparameters in the training process such as the learning rate and initialization variance, then derive principled choices for these hyperparameters to achieve some desired goal (e.g.\ improve feature learning). Examples of the infinite-width limit include works on initialization schemes such as \cite{he2016deep,yang2019scaling}, or more holistically network parametrizations such as \citep{yang2021tensor} where the authors introduced $\mu$P, a neural network parameterization ensuring feature learning in the infinite-width limit, offering precise scaling rules for architecture and learning rates to maximize feature learning. Examples for the depth limit include initialization strategies \citep{schoenholz2017deep, he2023deep, hayou19activation}, block scaling (see e.g. \cite{pmlr-v130-hayou21a,hayou2023on, noci2023shaped}), depth parametrizations \citep{yang2023depth,bordelon2023depthwise} etc. Here we propose to use the same strategy to derive scaling rules for the learning rate in LoRA for finetuning. More precisely, we study the infinite-width limit of LoRA finetuning dynamics and show that standard LoRA setup is suboptimal. We correct this by introducing a new method called LoRA$+$ that improves feature learning in low rank adaptation in the this limit. The key innovation in LoRA$+$ is setting different learning rates for $A$ and $B$ modules (LoRA modules) as explained in \cref{fig:intro_table}. Our theory is validated with extensive empirical results with different language of models and tasks.
